\centerline{\bf \huge 7 класс}
\centerline{\normalsize{\it Работа рассчитана на \bf{120} минут}}

\problem{7.1}
В лесу жили два гнома - Арсений и Вениамин. Арсений врал по понедельникам, вторникам и средам, а Вениамин - по вторникам, четвергам и субботам. В остальные дни они говорили только правду. Однажды при встрече одному из них задали вопрос: «Как тебя зовут?»

- Арсений, - ответил он.

- А какой сегодня день недели?

- Вчера было воскресенье.

Его приятель добавил: «А завтра будет пятница». У приятеля спросили: «Ты говоришь правду?».

- Я всегда говорю правду по средам, - ответил он.

Кто из них Арсений, кто - Вениамин и в какой день недели состоялся этот разговор?

\problem{7.2}
Малыш подарил Карлсону большую коробку конфет. Карлсон съел все конфеты за три дня. В первый день он съел $\dfrac{1}{5}$ всей коробки и еще 16 конфет. Во второй день $\dfrac{3}{10}$ остатка и еще 20 конфет. В третий день $\dfrac{3}{4}$ остатка и последние 30 конфет. Сколько конфет было в коробке?

%1-2 задачи Краснодар муницип 6 класс 2017 

\problem{7.3}
Найдите наименьшее натуральное число, которое при делении на $3,4, \ldots, 9$ даёт остатки $1, 2, \ldots, 7.$ Ответ обоснуйте.

\problem{7.4}
В треугольнике $A B C A B=B C$. Точка $E$ лежит на продолжении стороны $A C$ за точку $A$, причём $E A=A C$. Точка $F$ лежит на продолжении стороны $B C$ за точку $C$, причём $F C=C B$. Прямые $F A$ и $B E$ пересекаются в точке $D$. Докажите, что $A D=E D$.

%Тульская область 2024